% =====================================================================
% 前言 (Préface), 导论 (Introduction), 说明 (Avertissement)
% =====================================================================

\chapter*{前言 (Préface)}

本文是《平展覆盖与基本群》(Révêtements Étales et Groupe Fondamental)，Lecture Notes in Mathematics 224，Springer-Verlag，Berlin-Heidelberg-New York，1971，Alexander Grothendieck 等人一书的重排注释版。

LaTeX2e 的排版由一批志愿者完成，他们参与了一个由 Bas Edixhoven 主持的项目，有关详情见网址 http://www.math.leidenuniv.nl/\textasciitilde edix/. 版式由法国数学会 (Société mathématique de France) 最终完成.

这是对原文略作订正的一个版本. 它由法国数学会 (http://smf.emath.fr/) 在``数学文献''(Documents Mathématiques) 丛书中出版. Michel Raynaud 作了一些更新说明. 这些说明以方括号 [ ] 界定，并用记号 (MR) 标出 (第 211，220，256，292 页的注记以及第 67 页的注 7). 为统一记号，点 $x$ 的剩余域 (residue field) 记为 $k(x)$，局部环 $A$ 的剩余域记为 $k(A)$.

另有一个电子版本，意在复现原文.

两个版本共用一个源文件，该文件位于 arXiv.org 电子预印本服务器 http://arxiv.org/ 上. 两个版本之间的差异在源文件中均有记载.

旧页码已并入页边 (数字 $n$ 表示第 $n$ 页的开头).

\cleardoublepage

\chapter*{导论 (Introduction)}

在本导论的第一部分，我们对本卷的内容作一些说明；在第二部分，则对``Bois-Marie 代数几何讨论班''整体作说明，本卷构成其第一卷.

\section*{1.}

本卷在``Kronecker 式''的观点下，呈现了代数几何 (Géométrie Algébrique) 中基本群 (fundamental group) 理论的基础，该观点使人能够以同等的地位处理通常意义下的代数簇 (algebraic variety) 的情形，以及例如一个数域 (number field) 的整数环 (ring of integers) 的情形. 这一观点只有在概形 (scheme) 的语言中才能得到令人满意的表述，我们将自由地使用这种语言，以及 J. Dieudonné 与 A. Grothendieck 的《代数几何原理》(Éléments de Géométrie Algébrique，下文引用为 EGA) 前三章所阐述的主要结果. 研读本卷``Bois-Marie 代数几何讨论班''不需要代数几何的其他知识，因此对一位希望熟悉这些技巧的读者而言，它可作为现代代数几何技术的一个入门.

本书的讲 (Exposé) I 至 XI 是讨论班口头发言油印笔记的逐字复制，几乎未加改动，这些笔记曾由高等科学研究院 (Institut des Hautes Études Scientifiques)\footnote{以及随后的讨论班笔记. 由于这种分发方式长此以往被证明是不切实际且不充分的，今后所有的``Bois-Marie 代数几何讨论班''都将像本卷一样以书籍形式出版.} 分发. 我们只限于在原始文本中加入一些脚注，订正若干排印错误，并作了术语上的调整，其中``简单态射''(simple morphism) 一词已被``光滑态射''(smooth morphism) 所取代，后者不会引起同样的混淆.

讲 I 至 IV 呈现平展态射 (etale morphism) 与光滑态射的局部概念；它们几乎不用《原理》(Éléments) 第一章所阐述的概形语言\footnote{如今在《原理》第四章 §§17 和 18 中可以找到更完整的研究.}. 讲 V 呈现概形的基本群的公理化描述，即使在经典情形 (该概形化为一个域的谱 (spectrum)) 下也是有用的，在那里人们找到通常的 Galois 理论的一个非常便利的重述. 讲 VI 和 VIII 呈现下降 (descent) 理论，近年来下降理论在代数几何中获得了日益增长的重要性，并且它可能为解析几何 (Géométrie Analytique) 和拓扑学 (Topologie) 提供类似的服务. 应当指出，讲 VII 当时未曾写出，其内容已被并入 J. Giraud 的一篇著作 (Méthode de la Descente，Bull. Soc. Math. France，Mémoire 2，1964，viii + 150 p.). 在讲 IX 中，我们更专门地研究平展态射的下降，为基本群的 Van Kampen 型定理得到一种系统的方法，这些定理在这里表现为下降定理的简单翻译. 本质上，这是计算连通概形 $X$ 的基本群的一种方法，它给定了满射且紧合的 (proper) 态射，比方说 $X' \to X$，用 $X'$ 的连通分支的基本群、纤维积 $X' \times_X X'$，$X' \times_X X' \times_X X'$，以及上述概形之间的典范单纯 (simplicial) 态射在这些群之间诱导的同态来表达. 讲 X 给出基本群的特殊化 (specialization) 理论，针对紧合且光滑的态射，其最引人注目的结果在于 (几乎完全地) 确定特征 $p > 0$ 中光滑代数曲线的基本群，这得益于特征零中已知的超验方法的结果. 讲 XI 给出一些例子和补充，特别以同调的形式阐明了 Kummer 覆盖与 Artin-Schreier 覆盖的理论. 关于文本的其他评注，参见多印本版中的《说明》，它紧接本导论之后.

自从 1961 年本讨论班成文以来，M. Artin 和我本人合作发展了平展拓扑 (etale topology) 的语言以及相应的上同调理论，该理论在代数几何讨论班 SGA 4《概形的平展上同调》部分中阐述，即将在与本卷相同的丛书中出版. 这种语言以及它目前所拥有的成果，为基本群的研究提供了一个特别灵活的工具，使人们能够更好地理解并超越此处阐述的某些结果. 因此理应从这个观点出发彻底重做基本群理论 (事实上，所有关键结果目前都已出现在上引文献 (loc. cit.) 中). 这正是《原理》中拟用于基本群的那一章所计划的，它本应还包含若干无法在此占有一席之地的其他发展 (它们依据奇点解消 (resolution of singularities) 技术)：用该环的一个适当奇点解消来计算完备局部环的``局部基本群'' (local fundamental group)，无真性 (properness) 假设下基本群的局部与整体 Künneth 公式 (参见讲 XIII)，M. Artin 关于优的 (excellent) Hensel 局部环与其完备化的局部基本群之比较的结果 (SGA 4 XIX). 还应当指出，有必要发展一个意象 (topos) 的基本群理论，它将同时涵盖通常的拓扑理论、其半单纯 (semi-simplicial) 版本、本卷讲 V 中发展的``投射有限'' (profinite) 变体，以及 SGA 3 X 7 中稍更一般的投射离散 (pro-discrete) 变体 (它适用于非正规且非单支的概形的情形). 在从这一角度对理论进行全面改造之前，Raynaud 女士的讲 XIII 利用 SGA 4 的语言和结果，旨在表明在若干典型问题中可以得到的益处，特别是把讲 X 的某些结果推广到非紧合的相对概形. 其中特别给出了特征任意的非完备代数曲线的``素于 $p$''(prime to p) 基本群的结构 (这是我 1959 年宣布过的，但其证明至今未发表).

尽管有如此多的缺陷和不完善之处 (另一些人会说：正因为这些缺陷和不完善之处)，我认为本卷对于那些希望熟悉基本群理论的读者，以及作为一部参考著作，都是有用的，直至一部免于我刚列举的各种批评的文本写出并问世.

\section*{2.}

本卷构成``Bois-Marie 代数几何讨论班''的第 1 卷，其后续各卷预计在与本卷相同的丛书中出版. 讨论班的目标，与 J. Dieudonné 和 A. Grothendieck 的论著《代数几何原理》并驾齐驱，是按照后一部著作中的观点为代数几何奠定基础. 讨论班所有各卷的标准参考文献由《代数几何原理》第一、二、三章构成 (引用为 EGA I，II，III)，并假定读者已具备这些章所蕴含的交换代数 (commutative algebra) 与同调代数 (homological algebra) 的装备\footnote{关于这一点，参见 EGA I 的导论.}. 此外，讨论班的每一卷都将根据需要自由引用同一讨论班先前的各卷，或《原理》中已出版或即将出版的其他各章.

讨论班的每一部分都围绕一个主题，该主题在相应一卷或多卷的标题中给出；口头讨论班通常持续一个学年，有时更长. 讨论班每一部分内部的各讲之间通常具有密切的逻辑依存关系；相反，讨论班的不同部分在很大程度上彼此逻辑独立. 这样，``群概形'' (Schémas en Groupes) 部分几乎完全独立于在时间上先于它的讨论班的两个部分；相反，它频繁引用 EGA IV 的结果. 下面是讨论班即将出版的各部分的清单 (下文引用为 SGA 1 至 SGA 7)：

\begin{itemize}\tightlist
\item SGA 1. 平展覆盖与基本群 (Révêtements étales et groupe fondamental)，1960 和 1961.
\item SGA 2. 凝聚层的局部上同调与局部和整体 Lefschetz 定理 (Cohomologie locale des faisceaux cohérents et théorèmes de Lefschetz locaux et globaux)，1961/62.
\item SGA 3. 群概形 (Schémas en groupes)，1963 和 1964 (3 卷)，与 M. Demazure 合作.
\item SGA 4. 意象理论与概形的平展上同调 (Théorie des topos et cohomologie étale des schémas)，1963/64 (3 卷)，与 M. Artin 和 J.L. Verdier 合作.
\item SGA 5. $\ell$-adic 上同调与 $L$ 函数 (Cohomologie $\ell$-adique et fonctions $L$)，1964 和 1965 (2 卷).
\item SGA 6. 相交理论与 Riemann-Roch 定理 (Théorie des intersections et théorème de Riemann-Roch)，1966/67 (2 卷)，与 P. Berthelot 和 L. Illusie 合作.
\item SGA 7. 代数几何中的局部单值性群 (Groupes de monodromie locale en géométrie algébrique).
\end{itemize}

这些局部讨论班中有三个是与他人合作主持的，他们将作为合著者出现在相应各卷的封面上. 至于讨论班的其他活跃参与者，其作用 (无论在写作方面还是在数学定稿工作方面) 逐年增长，每位参与者的名字都列在他以演讲人或撰稿人身份负责的各讲之首，而某一特定卷中参与者的名单则在该卷的扉页上标明.

应当对讨论班与《原理》之间的关系作一些说明. 《原理》原则上旨在全面阐述在代数几何中被认为最基础的概念和技术，随着这些概念和技术本身在逻辑一致性与审美要求的自然作用之下显现出来. 从这个角度看，把讨论班视为《原理》的初版是自然的，它迟早几乎完全会被后者吸收. 这一进程在几年前就已经在一定程度上开始了，因为本卷 SGA 1 的讲 I 至 IV 已完全被 EGA IV 吸收，而讲 VI 至 VIII 本应在几年内被 EGA VI 吸收. 然而，随着《原理》和讨论班中建设的推进，以及整体规模日益明确，最初的宗旨 (按照它，讨论班不过是一个初步而临时的版本) 显得越来越不现实，原因 (除其他外) 在于深谋远虑的大自然为人的寿命设定的限度. 考虑到讨论班各部分的写作通常都很精心，也许只有在后来的写作进展能够在付出相当深刻的改动代价之下带来非常实质性的改进时，才有必要在《原理》(或接替它的论著) 中重写这样的部分. 当前的情形对 SGA 1 讨论班就是如此，如上所述，对 SGA 2 亦然 (得益于 Raynaud 女士近期的结果). 相反，目前没有任何迹象表明，上述 SGA 3 至 SGA 7 中任何一个部分在不久的将来会如此.

``Bois-Marie 代数几何讨论班''内部的参考文献是这样给出的. 对讨论班 SGA 1 至 SGA 7 某一部分的内部引用采用 III 9.7 的格式，其中数字 III 表示讲的编号，它出现在该讲的每一页的页首，而 9.7 是该讲内部陈述 (或定义、注记等) 的编号. 必要时可以使用更长的十进制数字，例如用 9.7.1，9.7.2 来指称命题 9.7 证明中的各个步骤. 引用 III 9 表示讲 III 的第 9 节. 对某一讲内部的引用省略讲的编号. 对讨论班其他部分的引用使用同样的记号，但前面冠以 SGA 相应部分的名称，如 SGA 1 III 9.7. 同样，引用 EGA IV 11.5.7 意指：《代数几何原理》第四章，陈述 (或定义等) 11.5.7；这里第一个阿拉伯数字仍表示节的编号. 除了在全部 SGA 中通行的这些约定外，与某一讲相关的文献 (bibliography) 通常汇集在该讲的末尾，讲内部按照惯例用方括号中的数字如 [3] 来引用.

最后，为读者方便起见，只要看起来有必要，我们将在 SGA 各卷的末尾附上一个记号索引 (index of notations) 和一个术语索引 (terminological index)，后者在适当处包含所用法语术语的英文翻译.

我要在本导论之后附上一段数学之外的评论. 1969 年 11 月，我得知高等科学研究院 (IHES) -- 我自其成立以来基本上是那里的教授 -- 三年来一直接受国防部 (Ministère des Armées) 的资助. 早在我作为初出茅庐的研究者时，我就发现大多数科学家愿意以这种或那种形式与军事机构合作，他们对此几乎毫无顾忌，这令我极为遗憾. 我当时的动机本质上属于道德性质，因此很难被认真对待. 今天，鉴于军事机构的泛滥及其所掌握的大规模毁灭手段对人类构成毁灭的危险，这些动机获得了新的力量与维度. 我已经在别处以更详细的方式说明过这些问题，它们比任何科学 (包括数学) 的进展都重要得多；例如可就此查阅 G. Edwards 在《Survivre》杂志第 1 期 (1970 年 8 月) 上的文章，它概括了我在别处就这些问题所作的更详尽的阐述. 就这样，我发现自己在一个机构里工作了三年，而它在我毫不知情的情况下参与了在我看来既不道德又危险的资助方式\footnote{不言而喻，我刚才表达的意见仅代表我个人的责任，而不代表出版本卷的 Springer 出版社的责任.}. 由于目前在 IHES 的同事中只有我持这种意见，这使我争取从 IHES 预算中取消军费资助的努力归于失败，于是我在 1970 年 9 月 30 日作出了势在必行的决定，离开了 IHES，并且只要该机构继续接受此类资助，我就暂停与该机构的一切科学合作. 我已请求 IHES 所长 Motchane 先生，使 IHES 从 1970 年 10 月 1 日起不再传播我署名的数学文本，或属于 Bois-Marie 代数几何讨论班的文本. 如上所述，本讨论班的传播将由 Julius Springer 出版社在 Lecture Notes 丛书中承担. 我在此高兴地感谢 Springer 出版社和 K. Peters 先生，感谢他们为出版本书所提供的有效而殷勤的帮助，特别是承担了新补入旧讨论班的各讲以及补全不全讨论班的各讲所需的照相胶印排版.

我还要感谢 J.P. Delale 先生，他承担了编制记号索引和术语索引这一吃力不讨好的工作.

\begin{flushright}
Massy，1970 年 8 月.
\end{flushright}

\cleardoublepage

\chapter*{说明 (Avertissement)}

已写成的每个讲都是若干连续口头发言内容的结晶. 似乎没有必要说明其具体日期.

讲 VII 在讲 VIII 中被反复引用，但它未曾由演讲人写出；在口头报告中，演讲人只限于从严格实用的观点勾勒一般范畴中下降的语言，而没有深入这种语言所引发的逻辑困难. 事实证明，对这种语言的正确阐述将超出本笔记的范围，仅就其篇幅而言也是如此. 关于下降理论的规范阐述，我请读者参阅 Jean Giraud 一篇正在准备中的文章. 在它问世之前\footnote{现已问世：J. Giraud，Méthodes de la descente，Mémoire no 2 de la Société mathématiques de France，1964.}，我认为细心的读者不难凭借自己的力量补足讲 VIII 中的``幻影引用'' (ghost references).

另有一些口头发言，位于讲 XI 之后，文本的某些地方曾提到它们，它们也未曾写出，本应构成一个讲 XII 和一个讲 XIII 的内容. 这些口头发言的前几个，在带幂零元 (nilpotent elements) 的概形与解析空间 (analytic space) 的框架内 (如 Cartan 讨论班 1960/61 所引入的)，重述了与一个在完备赋值域 (complete valued field) $k$ 上局部有限型的预概形 (prescheme) 相关联的解析空间的构造、$k$ 为复数域情形的 GAGA 型定理，以及它们在超验途径定义的基本群与本笔记所研究的基本群的比较上的应用 (比较 A. Grothendieck，Fondements de la Géométrie Algébrique，Séminaire Bourbaki no 190，第 10 页，1959 年 12 月). 后几个口头发言勾勒了本文所发展的方法向允许缓和分歧 (tame ramification) 的覆盖的研究、以及去掉有限个点后完备曲线的基本群结构的研究的推广 (比较上引文献 (loc. cit.)，no 182，第 27 页，定理 14). 这些发言没有引入任何本质上新的思想，因此在《代数几何原理》相应各章出版之前，似乎没有必要对它们作规范的书面整理\footnote{它们已包含在本卷中 Raynaud 女士的讲 XII 里，其证明与口头讨论班中所阐述的原始证明不同 (参见导论).}.

相反，针对基本群和 Picard 群的 Lefschetz 型定理，无论从局部还是整体观点看，都已构成 1962 年一个独立讨论班的对象，它已完全写出，可供使用者查阅\footnote{凝聚层的平展上同调与局部和整体 Lefschetz 定理 (SGA 2)，已由 North Holland Pub. Cie. 出版.}. 应当指出，无论是本讨论班还是 1962 年讨论班所发展的结果，都将在预概形平展上同调中若干关键结果的发表中得到本质性的使用，这些关键结果将构成一个 (由 M. Artin 和我本人主持的) 1963/64 年讨论班的对象，目前正在准备之中\footnote{概形的平展上同调 (引用为 SGA 4)，即将在同一丛书中出版.}.

讲 I 至 IV 本质上是局部性的、非常初等的，将被《代数几何原理》第四章完全吸收，该章的第一部分正在付印，大概将于 64 年底出版. 尽管如此，它们对于一位希望在深入系统论著的奥秘之前了解光滑、平展或平坦态射基本性质的读者仍是有用的. 至于其他各讲，它们将被《原理》第八章吸收\footnote{事实上，由于《原理》最初计划的改变，基本群的研究被推迟到比刚才所指出的章更靠后的某一章. 比较位于本《说明》之前的导论.}，而该章在数年之内很难考虑出版.

\begin{flushright}
Bures，1963 年 6 月.
\end{flushright}

\cleardoublepage
